% EXCERPT: original lines 1-25
\section{Gaussian scales}\label{sec:gaussian}

We next identify the fluctuation law on scales where the independent
common-word increments have negligible large jumps. There are three
steps: an annealed limit for independent paths, a comparison for shared
paths away from one another, and an occupation estimate excluding a
singular interaction on their projected diagonal. The resulting shared
two-copy limit will imply a quenched limit in probability.

\begin{definition}[Gaussian sequence]\label{gauss-sequence}
A sequence \(r\to\infty\) is \emph{Gaussian} if
\[
n(r)\Pr(|S|>a r)\longrightarrow0\qquad\text{for every fixed }a>0.
\tag{14}\label{eq:14}
\]
\end{definition}

Throughout this section limits are along any fixed Gaussian sequence.
Height-indexed processes are linearly interpolated on the grid
\(h/n(r)\). Functional convergence means weak convergence in
\(C([0,T],\mathbb R^k)\), with the uniform norm, for every fixed finite
\(T\); a slightly larger horizon is used when interpolating near its
endpoint. Put \(\sigma_*^2=1/(2\bar L)\).

\subsection{Independent-path Brownian limit}

% EXCERPT: original lines 370-428
\begin{lemma}[Off-diagonal comparison]\label{gauss-off-diagonal}
Fix \(a,t>0\). Uniformly over same-layer starts with
\(|y_2-x_2|\ge ar\), the total variation distance between the pair
path-prefix laws through level \(k=\lfloor tn(r)\rfloor\), under
\(P_{x,y}^{++}\) and under independent conditioned walks, satisfies
\[
\limsup {\rm dist}_{\rm TV}\le C\exp(-ca^2/t).
\tag{17}\label{eq:17}
\]
The constants depend only on the walk law.
\end{lemma}

\begin{proof}
In the top layer choose the windows
\[
I_1=\{z:|z_2-x_2-k\theta_r|<ar/3\},\qquad
I_2=\{z:|z_2-y_2-k\theta_r|<ar/3\}.
\]
They are disjoint. Lemma~\ref{gauss-shared-marginals} and the normal
tail bound show that the shared conditioned probability of either
endpoint error has limsup at most \(Ce^{-ca^2/t}\), uniformly in the
starts. Average the contact bound under the environment
law with density \(q(x)q(y)/p_2(x,y)\). This gives the same bound for
trace contact under \(P_{x,y}^{++}\).

To compare laws off contact, first use raw successful prefixes: both
walks reach level \(k\) before dropping below the common base. For any
two disjoint finite traces, their raw prefix weights agree after
averaging in shared and independent environments, because their
departure-row sets are disjoint. Passing from two finite successes to
two infinite no-drop events removes raw mass at most \(2(u_k-p)\)
under either pair law. Moreover \(p_2(x,y)-p^2\to0\) uniformly for the
present starts, since their separation tends to infinity. Thus the
conditioned prefix laws agree off contact up to a variation error that
tends to zero. Their total masses are both one, so the contact mass on
the independent side has the same upper bound up to this error. The
contact estimate proves \eqref{eq:17}.
\end{proof}

\subsection{Zero occupation on the projected diagonal}

The preceding comparison applies when the second-coordinate gap is of
order \(r\). It does not exclude a limiting pair that spends positive
time at zero gap. We now bound that occupation. The estimate counts
common true cuts first; exact-truth factorization will then transfer
it to every deterministic layer.

\begin{proposition}[Vanishing diagonal occupation]\label{gauss-diagonal}
Along every Gaussian sequence, under \(P_{0,0}^{++}\), for each fixed
finite \(T\),
\[
\lim_{\rho\downarrow0}\limsup_{r\to\infty}
\frac1{n(r)}E_{0,0}^{++}
 \sum_{0\le h\le Tn(r)}\mathbf1_{\{|Z_h|\le\rho r\}}=0.
\tag{18}\label{eq:18}
\]
\end{proposition}

\begin{proof}

% EXCERPT: original lines 654-657
\begin{proposition}[Independent joint limit for shared paths]\label{gauss-joint}
Under \(P_{0,0}^{++}\), the pair of processes in \eqref{eq:16}
converges functionally to two independent Brownian motions, each with
variance parameter \(\sigma_*^2\).

